	%iffalse
	\let\negmedspace\undefined
	\let\negthickspace\undefined
	\documentclass[journal,12pt,twocolumn]{IEEEtran}
	\usepackage{cite}
	\usepackage{amsmath,amssymb,amsfonts,amsthm}
	\usepackage{algorithmic}
	\usepackage{graphicx}
	\usepackage{textcomp}
	\usepackage{xcolor}
	\usepackage{txfonts}
	\usepackage{listings}
	\usepackage{enumitem}
	\usepackage{mathtools}
	\usepackage{gensymb}
	\usepackage{comment}
	\usepackage[breaklinks=true]{hyperref}
	\usepackage{tkz-euclide} 
	\usepackage{listings}
	\usepackage{gvv}                                        
	%\def\inputGnumericTable{}                                 
	\usepackage[latin1]{inputenc}                                
	\usepackage{color}                                            
	\usepackage{array}                                            
	\usepackage{longtable}                                       
	\usepackage{calc}                                             
	\usepackage{multirow}                                         
	\usepackage{hhline}                                           
	\usepackage{ifthen}                                           
	\usepackage{lscape}
	\usepackage{tabularx}
	\usepackage{array}
	\usepackage{float}


	\newtheorem{theorem}{Theorem}[section]
	\newtheorem{problem}{Problem}
	\newtheorem{proposition}{Proposition}[section]
	\newtheorem{lemma}{Lemma}[section]
	\newtheorem{corollary}[theorem]{Corollary}
	\newtheorem{example}{Example}[section]
	\newtheorem{definition}[problem]{Definition}
	\newcommand{\BEQA}{\begin{eqnarray}}
	\newcommand{\EEQA}{\end{eqnarray}}
	\newcommand{\define}{\stackrel{\triangle}{=}}
	\theoremstyle{remark}
	\newtheorem{rem}{Remark}

	% Marks the beginning of the document
	\begin{document}
	\bibliographystyle{IEEEtran}
	\vspace{3cm}

	\title{18. Definite Integrals And Applications of Integrals.}
	\author{AI24BTECH11016 - Jakkula Adishesh Balaji}
	\maketitle
	\newpage
	\bigskip

	\renewcommand{\thefigure}{\theenumi}
	\renewcommand{\thetable}{\theenumi}

	Section-B | JEE Main / AIEEE
	\begin{enumerate}
    \setcounter{enumi}{15}
	      \item 
		      The area of the region bounded by the curves $y=|x-2|, x=1, x=3$ and the x-axis is
		     \begin{enumerate}
		              \item 4
		              \item 2
		              \item 3
		              \item 1
		     \end{enumerate}
		     \hfill [2004] 
	      \item
           If $I_1=\int_{0}^{1} 2^{x^{2}} \,dx,I_2=\int_{0}^{1} 2^{x^{3}} \,dx,I_3=\int_{1}^{2} 2^{x^{2}} \,dx, I_4=\int_{1}^{2} 2^{x^{3}} \, dx
          $ then 
          \begin{enumerate}
                   \item $I_2>I_1$
                   \item $I_1>I_2$
                   \item $I_3=I_4$
                   \item $I_3>I_4$
          \end{enumerate}
          \hfill [2005]
	      \item
		      The area enclosed between the curve $y=log_e(x+e)$ and the coordinate  axes is 
		     \begin{enumerate}
		              \item 1
		              \item 2
		              \item 3
		              \item 4
		     \end{enumerate}
		     \hfill [2005] 
	      \item 
		      The parabolas $y^{2}=4x$ and $x^{2}=4y$ divide the square region bounded by the lines $x=4, y=4$ and the coordinate axes. If $S_1, S_2, S_3$ are respectively the areas of these parts numbered from top to bottom then $S_1:S_2:S_3$ is 
		     \begin{enumerate}
		             \item 1:2:1
		             \item 1:2:3
		             \item 2:1:2
		             \item 1:1:1
		     \end{enumerate}
		     \hfill [2005] 
	     \item 
		      Let f(x) be a non-negative continuous function such that the area bounded by the curve $y=f(x)$, x-axis and the oordinates $x=\pi/4$ and $x=\beta>\pi/4$ is $((\beta sin(\beta)+\pi/4cos(\beta))$. Then $f(\pi/2)$ is
		     \begin{enumerate}
		            \item $\frac{\pi}{4}+\sqrt{2}-1$
		            \item $\frac{\pi}{2}-\sqrt{2}+1$
		            \item $1-\frac{\pi}{4}-\sqrt{2}$
		            \item $1-\frac{\pi}{4}+\sqrt{2}$
		     \end{enumerate}
		     \hfill [2005] 
	     \item 
		      The value of $\int_{-\pi}^{\pi} \frac{cos^{2}(x)}{1+a^{x}}\,dx$, $a>0$, is
		     \begin{enumerate}
		           \item $\pi$
		           \item $\frac{\pi}{2}$
		           \item $\frac{\pi}{a}$
		           \item $2\pi$
		     \end{enumerate}
		     \hfill [2005] 
	     \item 
		     The value of the integral $\int_{3}^{6} \frac{\sqrt{x}}{\sqrt{9-x}+\sqrt{x}} \,dx$ is
		    \begin{enumerate}
		          \item $\frac{1}{2}$
		          \item $\frac{3}{2}$
		          \item 2
		          \item 1
		    \end{enumerate}
		    \hfill [2006] 
	     \item 
		     $\int_{0}^{\pi} xf(sinx) \,dx$ is equal to
		    \begin{enumerate}
		         \item $\pi (\int_{0}^{\pi} f(cosx) \,dx)$
		         \item $\pi (\int_{0}^{\pi} f(sinx) \,dx)$
		         \item $\frac{\pi}{2} (\int_{0}^{\frac{\pi}{2}} f(sinx) \,dx)$
		         \item $\pi (\int_{0}^{\frac{\pi}{2}} f(cosx) \,dx)$
		    \end{enumerate}
		    \hfill [2006] 
	     \item
		    $\int_{\frac{-3\pi}{2}}^{\frac{-\pi}{2}} [(x+\pi)^{3}+cos^{2}(x+3\pi)] \,dx$ is equal to
		   \begin{enumerate}
		        \item $\frac{\pi^{4}}{32}$
		        \item $\frac{\pi^{4}}{32}+\frac{\pi}{2}$
		        \item $\frac{\pi}{2}$
		        \item $\frac{\pi}{4}-1$
		   \end{enumerate}
		   \hfill [2006] 
	    \item 
		    The value of $\int_{1}^{a} [x]f^{'}(x) \,dx$, $a>1$ where [x] denotes the greatest integer not exceeding x is
		   \begin{enumerate}
		       \item $af(a)-{f(1)+f(2)+...f([a])}$
		       \item $[a]f(a)-{f(1)+f(2)+...f([a])}$
		       \item $[a]f([a])-{f(1)+f(2)+...f(a)}$
		       \item $af([a])-{f(1)+f(2)+...f(a)}$
		   \end{enumerate}
		   \hfill [2006] 
	   \item 
		    Let $F(x)=f(x)+f(\frac{1}{x})$, where $f(x)=\int_{1}^{x} \frac{logt}{1+t} \,dt$, Then $F(e)$ equals
		   \begin{enumerate}
		      \item 1
		      \item 2
		      \item $\frac{1}{2}$
		      \item 0
		   \end{enumerate}
		   \hfill [2007] 
	   \item 
    		    The solution for x of the equation $\int_{\sqrt2}^{x} \frac{1}{t\sqrt{t^{2}-1}} \,dt$ is
		   \begin{enumerate}
		      \item $\frac{\sqrt{3}}{2}$
		      \item $2\sqrt{2}$
		      \item 2
		      \item none
		   \end{enumerate}
		   \hfill [2007] 
	   \item 
		  The area enclosed between the curves $y^{2}=x$ and $y=|x|$ is	
		 \begin{enumerate}
		     \item $\frac{1}{6}$
		     \item $\frac{1}{3}$
		     \item $\frac{2}{3}$
		     \item 1
		 \end{enumerate}
		 \hfill [2007] 
	   \item 
		 Let $I=\int_{0}^{1} \frac{sinx}{\sqrt{x}} \,dx$ and $J=\int_{0}^{1} \frac{cosx}{\sqrt{x}} \,dx$. Then which one of the following is true?
		\begin{enumerate}
		    \item $I>\frac{2}{3}$ and $J>2$
		    \item $I<\frac{2}{3}$ and $J<2$
		    \item $I<\frac{2}{3}$ and $J>2$
		    \item $I>\frac{2}{3}$ and $J<2$
		\end{enumerate}
		\hfill [2007] 
	   \item 
		 The area of the plane region bounded by the curves $x+2y^{2}=0$ and $x+3y^{2}=1$ is equal to 
		\begin{enumerate}
		   \item $\frac{5}{3}$
		   \item $\frac{1}{3}$
		   \item $\frac{2}{3}$
		   \item $\frac{4}{3}$
		\end{enumerate}
		\hfill [2008] 
	\end{enumerate}
	\end{document}
